\documentclass[11pt]{article}
\usepackage[T1]{fontenc}
\usepackage{lmodern}
\usepackage[letterpaper,margin=0.75in]{geometry}
\usepackage{amsmath,amssymb,microtype,xcolor,hyperref}
\hypersetup{colorlinks=true,urlcolor=blue!45!black,linkcolor=blue!45!black,
 filecolor=blue!45!black,
 pdftitle={Assessment of the OpenAI November 4 Yang-Mills submission},
 pdfauthor={Yang--Mills group -- working draft}}
\setlength{\parindent}{0pt}
\setlength{\parskip}{5pt}
\setlength{\abovedisplayskip}{6pt}
\setlength{\belowdisplayskip}{6pt}
\pagestyle{empty}
\clubpenalty=10000
\widowpenalty=10000
\newcommand{\reviewdoc}[2]{\href{https://github.com/yang-mills-group/submissions/blob/main/reviews/openai-20261104/#1}{#2}}

% Separate draft: the October assessment is preserved unchanged.
% Repository links currently require access to the private review repository.
% Supporting notes are published under docs/matching/.
\begin{document}
\begin{center}
{\Large\bfseries OpenAI Yang--Mills: November assessment}\par
\vspace{3pt}
\href{https://github.com/yang-mills-group/OpenAI-YM-20261104/tree/c1bb5a5566b2bad9d593980d470439f815895317/}{4 November 2026 submission}\quad\textbullet\quad Yang--Mills group\par
{\small Draft for discussion, 16 November 2026}
\end{center}
\vspace{-4pt}

This supplements our \href{https://github.com/yang-mills-group/YM/blob/main/docs/assessment-1.pdf}{October assessment}.
November adds all-orders matching and passes our formal checks.
Human review remains minimal; a more extensive review is underway.

\textbf{1. Existence and completed checks.}
Both releases construct a Wightman theory from four-dimensional Wilson lattice
gauge theory in the stated simple-group setting, via a full fixed-volume
ultraviolet limit and a subsequential infinite-volume limit. November clarifies
the lattice-field dictionary and normalization, strengthens the phase-uniform
free-Maxwell comparison, and adds OPE consistency.
Our October checks transferred to the same November witness: Wilson origin,
OS axioms with growth on diagonal-flat joint Schwartz tests, Ward and stress
identities, field coverage and OPE bounds. Both combined targets passed
Comparator with standard Lean axioms only. Neither release establishes a mass
gap or uniqueness of the infinite-volume limit.

\textbf{2. Perturbative matching: statement and evidence.}
One Wilson construction and perturbative reference are chosen before the
volume sequence. Each selected continuum hierarchy has one calibrated field
frame $J_r$ and measured coupling $g_X(r)\to0$, common to all orders. Set
$S_{\sigma,r}(v;x,o)=S_\sigma(J_rv_1,\ldots,J_rv_n;o+rx_1,\ldots,o+rx_n)$
for local gauge-invariant fields. For every finite order $N$, derivative order
$j$, field-degree cutoffs and compact set $K$ of separated shapes,
\[
 \left\|D_x^\ell\!\left[S_{\sigma,r}(v;x,o)
       -\sum_{k=0}^{N}g_X(r)^k C^X_{n,k}(v;x)\right]\right\|
 \le A\!\prod_i\|v_i\|_{\mathrm{coeff}}\,|g_X(r)|^{N+1},
 \qquad \ell\le j.
\]
for sufficiently small $r>0$, uniformly in $x\in K$, center, phase and fields.
The constant and radius threshold may depend on $N,j,K$ and the cutoffs.
Analogous bounds control the full OPE output at each finite output cutoff.
The $C^X$ come from the submitted graph reference, intended as Landau-gauge
$\overline{\mathrm{MS}}$ with a fixed field/coupling conversion. The proof compares the
continuum limit of lattice coefficients with continuum flow solutions,
matching sums at each order. It does not assert convergence of the series.

The supplied Existence and Matching targets passed Comparator; our Lean
restatement checks the quantifiers and common witnesses. Independent finite
ghost- and gluon-bubble identities also passed Lean and Comparator, including
the finite $1/3$ polarization-trace term. The four-leg ghost box has analytic,
structural Lean and numerical checks; its final evaluator equality remains
unproved in Lean.

\textbf{3. Subtleties still requiring comparison at all orders.}
\begin{list}{\textbullet}{\setlength{\leftmargin}{12pt}\setlength{\itemsep}{2pt}\setlength{\parsep}{0pt}\setlength{\topsep}{2pt}}
\item \textbf{Renormalization.} Compare recursive scale-sector counterterms with
conventional forest subtraction, including overlapping subdivergences; fix
finite terms, the $\overline{\mathrm{MS}}$ scale and conversion to the measured scheme.
\item \textbf{Dimensional continuation.} Keep open tensor indices and evanescent terms
through pole subtraction: they can leave finite contributions. Check the order
of subtraction and infrared-cutoff removal.
\item \textbf{Composite operators.} Specify interacting completions of free
curvature labels, operator mixing, equation-of-motion and BRST-exact terms,
evanescent lifts and normalization. The label $F_0^2$ alone does not identify
$[F^2]_{\overline{\mathrm{MS}}}$.
\item \textbf{Flow and uniqueness.} Identify graph multiplicities, graded signs,
classical input and boundary data before using uniqueness to identify the
intended perturbative coefficients.
\item \textbf{OPE and scope.} Check the common field conversion and its inverse on
OPE outputs, compatibility across cutoffs and control against arbitrary exterior
insertions. Separated-point bounds do not establish contact identities or an
off-shell operator algebra.
\end{list}


{\small\textbf{Supporting working material (AI-generated):}
\reviewdoc{docs/existence/october-november-existence-comparison.pdf}{Existence comparison};
\reviewdoc{docs/matching/all-orders-statement-review.md}{matching statement};
\reviewdoc{docs/matching/november-matching-assessment-draft.pdf}{detailed matching review};
\reviewdoc{docs/matching/ghost-finite-check.md}{ghost check};
\reviewdoc{docs/matching/gluon-bubble-finite-identity.pdf}{gluon calculation}.
Repository links require group access.}
\end{document}
